\section{Notation and previous results}\label{sec:preliminaries}

The set of all natural numbers is identified with the first infinite ordinal, $\omega$.
$\Sym$ denotes the group of all permutations of $\omega$.
We identify a real series with a sequence
$a=\langle a_n:n<\omega\rangle\in\mathbb R^\omega$ and write
\[
 S_m(a)=\sum_{n<m}a_n,
 \qquad
 S_m^\pi(a)=\sum_{n<m}a_{\pi(n)}
 \quad(\pi\in\Sym, m \in \omega).
\]
A series $a$ is conditionally convergent if $\sum_n a_n$ converges but
$\sum_n\abs{a_n}$ diverges.
The set of all conditionally convergent series is denoted by $\CC$.

\begin{definition}\label{def:rr}
A collection $\Pi\subseteq\Sym$ is \emph{rearranging} if for every
$a\in\CC$ there exists $\pi\in\Pi$ such that
$\sum_n a_{\pi(n)}$ either diverges or converges to a real number different
from the original sum.  The \emph{rearrangement number} is defined as
\[
      \rr=\min\{\abs{\Pi}:\Pi\subseteq\Sym
      \text{ is rearranging}\}.
\]
\end{definition}


For
$f,g\in\omega^\omega$, write $f\leq^*g$ when
$f(n)\leq g(n)$ for all but finitely many $n\in \omega$.
If $f\not \leq^*g$, we write $g<^\infty f$.


The bounding number
$\bb$ is the least cardinality of a $\leq^*$-unbounded subset of
$\omega^\omega$.

Let $\cM$ and $\cN$ be the ideals of meager and Lebesgue null subsets
of $\mathbb R$, respectively, and let $\mathfrak c=|\mathbb R|$.
The cardinal $\covN$ is the least cardinality of a collection of Lebesgue null subsets of $\mathbb R$ whose union is $\mathbb R$.
Finally, the cardinal $\nonM$ is the least cardinality of a nonmeager subset of $\mathbb R$.
It is well known that $\aleph_1\leq \max\{\covN,\bb\}\leq\nonM\leq \mathfrak c$.
See \cite{BlassHandbook} for more
background on these cardinal invariants.

\begin{theorem}[{\cite[Theorems 7, 8, and 11]{BBBHHL}}]
\label{thm:known-sandwich}
In ZFC, $\max\{\covN,\bb\}\leq\rr\leq\nonM$.
\end{theorem}
For a sequence $r:\omega\to\omega\setminus\{0\}$, let
\[
 {\mathcal S}_r=\left\{\varphi:\omega\to[\omega]^{<\omega}:
                         \abs{\varphi(n)}\leq r(n)\text{ for every }n\right\}.
\]
The elements of ${\mathcal S}_r$ are called \emph{$r$-slaloms}. A function
$x\in\omega^\omega$ \emph{eventually avoids} $\varphi$ if
$x(n)\notin\varphi(n)$ for all but finitely many $n$.

\begin{theorem}[Bartoszy\'nski, {\cite[Lemma 2.4.8]{BartoszynskiJudah}}]\label{thm:slaloms}
For every $r:\omega\to\omega\setminus\{0\}$, every family
$\Phi\subseteq{\mathcal S}_r$ of cardinality less than $\nonM$
has a common eventual avoider.
\end{theorem}
We identify a natural number $m$ with $\{0,\ldots,m-1\}$. For a finite set $X$, $\operatorname{Sym}(X)$
denotes its permutation group.
In particular, for a natural number $L$, $\operatorname{Sym}(L)$ is the standard permutation group $S_L$.

Following \cite{BlassHandbook}, we use the language of Galois--Tukey connections.

\begin{definition}
      A \emph{relation triple}, or simply a \emph{relation}, is a triple $\mathbf A=(A_-,A_+,A)$, where $A_-$ and $A_+$ are nonempty sets and $A\subseteq A_-\times A_+$ has domain $A_-$.

      Elements of $A_-$ are called \emph{challenges} and elements of $A_+$ are called \emph{responses}.

      The norm of a relation $\mathbf A$ is defined as
      \[
            \|\mathbf A\|=\min\{|Y|:Y\subseteq A_+\text{ and }\forall x\in A_-\ \exists y\in Y\ (x,y)\in A\}.
      \]
\end{definition}

Thus:
\begin{itemize}
      \item The cardinal $\mathfrak b$ is $\|\mathfrak B\|$, where $\mathfrak B=(\omega^\omega,\omega^\omega,<^{\infty})$.
      \item The cardinal $\nonM$ is $\|\mathfrak S_r\|$, where $\mathfrak S_r=(\omega^\omega,{\mathcal S}_r,\in^\infty)$ for any $r:\omega\to\omega\setminus\{0\}$ with $\lim_{n\to\infty}r(n)=\infty$.
      Here $e\in^\infty\varphi$ means that $e(n)\in\varphi(n)$ for infinitely many $n$.
      \item The cardinal $\rr$ is $\|\mathfrak R\|$, where $\mathfrak R=(\CC,\Sym,\text{rearranges})$.
\end{itemize}

\begin{definition}
      Let $\mathbf A=(A_-,A_+,A)$ and $\mathbf B=(B_-,B_+,B)$ be relations. A \emph{morphism} from $\mathbf A$ to $\mathbf B$ is a pair of functions $(\phi_-,\phi_+)$, where $\phi_-:B_-\to A_-$ and $\phi_+:A_+\to B_+$, such that
      \[
            \forall x\in B_-\ \forall y\in A_+\ (\phi_-(x),y)\in A\rightarrow(x,\phi_+(y))\in B.
      \]
\end{definition}

It is straightforward to check that if there is a morphism from $\mathbf A$ to $\mathbf B$, then $\|\mathbf A\|\geq\|\mathbf B\|$.

\begin{definition}[{\cite[Definition 4.10]{BlassHandbook}}]
\label{def:sequential-composition}
For relations $\mathbf A=(A_-,A_+,A)$ and $\mathbf B=(B_-,B_+,B)$,
their \emph{sequential composition} is
\[
 \mathbf A;\mathbf B
 =\bigl(A_-\times{}^{A_+}B_-,\ A_+\times B_+,\ C\bigr),
\]
where ${}^{A_+}B_-$ is the set of functions from $A_+$ to $B_-$, and
\[
 (x,f)\,C\,(a,b)
 \quad\Longleftrightarrow\quad x\,A\,a\text{ and }f(a)\,B\,b.
\]
Thus the second challenge $f(a)$ depends on the response $a$ to the
first challenge $x$.
\end{definition}

\begin{remark}[{\cite[Theorem 4.11(3)]{BlassHandbook}}]
\label{rem:sequential-norm}
The norm of the sequential composition satisfies
\[
 \|\mathbf A;\mathbf B\|=\|\mathbf A\|\cdot\|\mathbf B\|.
\]
For the upper bound, which is all we need, take families $Y\subseteq A_+$
and $Z\subseteq B_+$ witnessing the two norms. Given a challenge $(x,f)$,
choose $a\in Y$ with $x\,A\,a$, and then $b\in Z$ with $f(a)\,B\,b$.
Hence $Y\times Z$ responds to every challenge in $\mathbf A;\mathbf B$.
\end{remark}
